% BHU PYQ -- Manually transcribed & error-corrected
\documentclass[11pt,a4paper]{article}

\usepackage[a4paper,top=2.5cm,bottom=2.5cm,left=2.8cm,right=2.8cm,headheight=14pt]{geometry}
\usepackage{amsmath,amssymb}
\usepackage[T1]{fontenc}
\usepackage[utf8]{inputenc}
\usepackage{lmodern}
\usepackage{microtype}
\usepackage{fancyhdr}
\usepackage{enumitem}
\usepackage{hyperref}
\usepackage{lastpage}
\usepackage{parskip}

\hypersetup{colorlinks=true,urlcolor=black,linkcolor=black,
  pdftitle={CHB-02A: Ancillary Chemistry-I -- BHU PYQ},pdfauthor={Banaras Hindu University}}

\pagestyle{fancy}\fancyhf{}
\renewcommand{\headrulewidth}{0.4pt}\renewcommand{\footrulewidth}{0.4pt}
\fancyhead[L]{\small\textit{Banaras Hindu University}}
\fancyhead[R]{\small\textit{CHB-02A: Ancillary Chemistry-I}}
\fancyfoot[C]{\small Page \thepage\ of \pageref{LastPage}}

\newlist{parts}{enumerate}{2}
\setlist[parts,1]{label=(\alph*),leftmargin=*,itemsep=5pt,topsep=3pt}
\setlist[parts,2]{label=(\roman*),leftmargin=*,itemsep=3pt,topsep=2pt}
\newcommand{\pts}[1]{\hfill\textit{[#1]}}

\begin{document}

\begin{center}
  {\large\bfseries BANARAS HINDU UNIVERSITY}\\[2pt]
  {\normalsize B.Sc. (Hons.) Semester II Examination, 2023-24}\\[6pt]
  \rule{0.8\linewidth}{0.5pt}\\[6pt]
  {\large\bfseries Paper No. CHB-02A: Ancillary Chemistry-I}\\[4pt]
  {\normalsize Time: 3 Hours\hfill Maximum Marks: 70}
\end{center}
\rule{\linewidth}{0.4pt}
\small
\noindent\textit{Instructions: Answer five questions in total, including Question~1 which is compulsory.}
\normalsize
\bigskip

\noindent\textbf{Question 1.} Answer all parts of the following: \pts{5 $\times$ 2 = 10}
\begin{parts}
  \item What is glycogen? Where is it stored in the human body?
  \item What do you mean by expanding octet and incomplete octet? Explain with one example of each.
  \item In which compounds can the hydrogen atom have the oxidation number $-1$? Explain with two examples.
  \item Explain the changes in the oxidation numbers of the atoms in the following reaction:
    \[ 2\text{KMnO}_4 + 10\text{FeSO}_4 + 8\text{H}_2\text{SO}_4 \to \text{K}_2\text{SO}_4 + 2\text{MnSO}_4 + 5\text{Fe}_2(\text{SO}_4)_3 + 8\text{H}_2\text{O} \]
  \item Differentiate between binding energy and packing fraction.
\end{parts}

\medskip\rule{\linewidth}{0.2pt}

\noindent\textbf{Question 2.}
\begin{parts}
  \item What is a spontaneous reaction? Derive the expression for the entropy change ($\Delta S$) of an ideal gas expanding isothermally. \pts{7}
  \item What is the van 't Hoff reaction isotherm? Derive the Gibbs-Helmholtz equation. \pts{7}
\end{parts}

\medskip\rule{\linewidth}{0.2pt}

\noindent\textbf{Question 3.}
\begin{parts}
  \item Differentiate between thin-layer chromatography (TLC) and paper chromatography. \pts{6}
  \item Describe the analytical applications of the following: (i) Dimethylglyoxime (DMG) and (ii) 8-hydroxyquinoline (oxine). \pts{8}
\end{parts}

\medskip\rule{\linewidth}{0.2pt}

\noindent\textbf{Question 4.}
\begin{parts}
  \item What is Fehling's solution? What are the analytical uses of Fehling's solution? \pts{4}
  \item How would you prepare Benedict's reagent? What happens when Benedict's reagent is treated with glucose in the presence of $\text{KSCN}$? \pts{6}
  \item Discuss the roles of the following metals in our body: (i) Cobalt and (ii) Zinc. \pts{4}
\end{parts}

\medskip\rule{\linewidth}{0.2pt}

\noindent\textbf{Question 5.}
\begin{parts}
  \item Describe the synthesis and uses of the following polymers: (i) Teflon (tetrafluoroethylene) and (ii) Buna-S rubber. \pts{6}
  \item Describe how iron is extracted from its haematite ore. Give all chemical reactions involved in the blast furnace. \dots \pts{6}
  \item Explain the discrepancy of bond angles in $\text{SO}_2$ and $\text{H}_2\text{O}$ using VSEPR theory. \pts{2}
\end{parts}

\vfill
\rule{\linewidth}{0.4pt}
\begin{center}\small
  \href{https://bhu-exam.vercel.app/}{Click here to visit our website}\\[4pt]
  \href{https://me-aryan.vercel.app/}{Click here to visit our contributor}\\[4pt]
  \href{https://quashan-me.vercel.app/}{Click here to visit our contributor}
\end{center}

\end{document}
